\chapter{Product Planning}

\section{Two-Session Overview}

Product planning converts the earlier opportunity work into a clear development direction. In these two sessions, your group will study the market, identify relevant technologies and constraints, and then convert that analysis into a mission statement that can guide later work on customer needs and product specifications.

By the end of Sessions 5 and 6, your group should be able to:
\begin{enumerate}
    \item identify relevant market segments, stakeholders, technologies, and constraints for the selected hydroponic product;
    \item explain how technical and physics-related considerations influence the planning decisions;
    \item draft and refine a mission statement for the project; and
    \item produce a shared planning document that will support Sessions 7 and 8.
\end{enumerate}

\section{Key Terms}

The following terms are used throughout this chapter:
\begin{itemize}
    \item \textbf{Market segment}: a group of customers with similar needs, behaviours, or use contexts.
    \item \textbf{Technology integration}: the choice of technical features or subsystems that could strengthen the product.
    \item \textbf{Mission statement}: a concise planning statement describing the product, its benefits, target market, and key goals.
    \item \textbf{Stakeholders}: people or groups affected by the product, including users, manufacturers, maintainers, sellers, and investors.
    \item \textbf{Assumptions and constraints}: the working conditions, limits, and uncertainties that shape the project.
\end{itemize}

\section{Why Product Planning Matters}

A good plan prevents later sessions from becoming disconnected. For example, if your team expects the product to be affordable for small urban growers, that assumption will affect customer needs, material choices, manufacturing approach, and target price. Similarly, if the concept depends on sensors or automated flow control, the team must already recognise the implications for calibration, power use, durability, and maintenance.

\section{Session 5: Individual Market Analysis and Draft Mission Statement}

Session 5 focuses on preparation. Each member should review the opportunity from earlier sessions and gather evidence that will help the team plan realistically.

\subsection{Activity 1: Market Analysis}

Before class, analyse the planning questions below for your selected product.

\begin{enumerate}
    \item \textbf{Market segments}: Which customer groups could use the product? For example, are you designing for household hobby growers, educational users, or commercial operators?
    \item \textbf{Technology opportunities}: Which technologies could improve the product? Examples include sensors, modular structures, dosing systems, or low-power control units.
    \item \textbf{Manufacturing and service considerations}: What production, maintenance, or service constraints might influence the product?
    \item \textbf{Target market and price}: Who is the main customer, and what price range is realistic in Malaysian Ringgit (RM)?
    \item \textbf{Assumptions, constraints, and stakeholders}: What conditions are you assuming, what limits exist, and who must be considered during planning?
\end{enumerate}

For Industrial Physics students, this analysis should also consider physical behaviour where relevant. For example, flow stability, corrosion resistance, thermal exposure, or sensor accuracy may affect both technology selection and product feasibility.

\subsection{Market Analysis Template}

Use Table~\ref{tab:session56-market-analysis} to organise your findings.

\begin{table}[htbp]
\centering
\caption{Example market analysis template for hydroponic product planning}
\label{tab:session56-market-analysis}
\renewcommand{\arraystretch}{1.2}
\begin{tabularx}{\textwidth}{|>{\raggedright\arraybackslash}p{0.17\textwidth}|>{\raggedright\arraybackslash}p{0.18\textwidth}|>{\raggedright\arraybackslash}X|>{\raggedright\arraybackslash}p{0.16\textwidth}|}
\hline
\textbf{Planning element} & \textbf{Example entry} & \textbf{Questions to answer} & \textbf{Sources or notes} \\
\hline
Market segment & Urban home grower & Who uses the product and what problem do they face? &  \\
\hline
Technology opportunity & Low-cost water-level sensor & What technology could improve function, safety, or usability? &  \\
\hline
Manufacturing or service issue & Easy cleaning of wetted parts & What production or maintenance constraint matters most? &  \\
\hline
Target market and price & Beginner hobby growers, RM150--RM250 & What price range is realistic and why? &  \\
\hline
Assumptions and constraints & Indoor use, limited power, small footprint & What limits must the concept respect? &  \\
\hline
Stakeholders & User, retailer, fabricator, maintainer & Who affects or is affected by the product? &  \\
\hline
\end{tabularx}
\end{table}

\subsection{Activity 2: Draft the Mission Statement}

After completing the analysis, write a short draft mission statement. A useful mission statement should include:
\begin{itemize}
    \item the product description;
    \item the main benefit proposition;
    \item the primary market and any secondary market;
    \item the key business or project goal; and
    \item the most important assumptions, constraints, and stakeholders.
\end{itemize}

\subsection{Mission Statement Template}

You may draft the mission statement using the structure in Table~\ref{tab:session56-mission-template} before converting it into a short paragraph.

\begin{table}[htbp]
\centering
\caption{Mission statement drafting template}
\label{tab:session56-mission-template}
\renewcommand{\arraystretch}{1.2}
\begin{tabularx}{\textwidth}{|>{\raggedright\arraybackslash}p{0.22\textwidth}|X|}
\hline
\textbf{Field} & \textbf{Draft content} \\
\hline
Product description &  \\
\hline
Benefit proposition &  \\
\hline
Primary market &  \\
\hline
Secondary market &  \\
\hline
Key goal &  \\
\hline
Assumptions and constraints &  \\
\hline
Stakeholders &  \\
\hline
\end{tabularx}
\end{table}

\subsection{Expected Session 5 Output}

Before the next class, each member should bring:
\begin{itemize}
    \item a completed market analysis table;
    \item one draft mission statement; and
    \item brief notes on supporting sources or evidence.
\end{itemize}

\section{Session 6: Consolidating the Group Plan}

Session 6 focuses on synthesis. Your group will compare the individual analyses, resolve differences, and produce one shared planning direction.

\subsection{Activity 1: Merge the Market Analyses}

Work through the following steps as a team:
\begin{enumerate}
    \item compare the market segments identified by each member;
    \item agree on the most relevant customer group or groups;
    \item discuss which technologies are attractive and technically realistic;
    \item review manufacturing, service, and maintenance issues that the product must respect;
    \item confirm a realistic target market and price range; and
    \item finalise the most important assumptions, constraints, and stakeholders.
\end{enumerate}

If opinions differ, use the innovation charter, user value, and project feasibility as the main criteria for deciding.

\subsection{Group Consolidation Template}

Use Table~\ref{tab:session56-group-consolidation} to capture the final consensus.

\begin{table}[htbp]
\centering
\caption{Group market analysis consolidation template}
\label{tab:session56-group-consolidation}
\renewcommand{\arraystretch}{1.2}
\begin{tabularx}{\textwidth}{|>{\raggedright\arraybackslash}p{0.17\textwidth}|>{\raggedright\arraybackslash}X|>{\raggedright\arraybackslash}p{0.18\textwidth}|}
\hline
\textbf{Planning element} & \textbf{Consensus decision} & \textbf{Rationale or evidence} \\
\hline
Market segment &  &  \\
\hline
Technology opportunity &  &  \\
\hline
Manufacturing and service &  &  \\
\hline
Target market and price &  &  \\
\hline
Assumptions and constraints &  &  \\
\hline
Stakeholders &  &  \\
\hline
\end{tabularx}
\end{table}

\subsection{Activity 2: Finalise the Mission Statement}

Once the group analysis is consolidated, convert it into one final mission statement. Check the statement against the following questions:
\begin{enumerate}
    \item Is the product described clearly?
    \item Is the value proposition meaningful for the selected market?
    \item Are the main assumptions and constraints realistic?
    \item Does the statement align with the innovation charter and shortlisted opportunity?
    \item Does the wording support later work on customer needs and specifications?
\end{enumerate}

\subsection{Collaboration Guidance}

To keep the discussion fair and efficient:
\begin{itemize}
    \item let each member present their strongest evidence briefly;
    \item record disagreements before voting or choosing a final direction;
    \item verify claims with credible sources where possible; and
    \item make sure the final wording is understood by the full team.
\end{itemize}

\subsection*{Optional Remarks (Extra Credit)}

If your lecturer allows extra credit, you may strengthen the submission by adding deeper justification. Useful remarks include:
\begin{itemize}
    \item market or competitor observations supported by credible sources;
    \item discussion of relevant trends or risks;
    \item justification for the selected target price; and
    \item explanation of why certain stakeholders were prioritised.
\end{itemize}

Use a consistent citation style such as APA if external sources are referenced.

\section{Summary and Next Steps}

At the end of Session 6, your group should have one refined market analysis and one final mission statement. These outputs become the planning foundation for Session 7, where you will identify customer needs, and Session 8, where those needs will be translated into measurable specifications.
